%%%%%%%%%%  scan idx 36  |  folio 23  %%%%%%%%%%
les
coefficients $c_n(\lambda^i(f))=c_n\text{ de }\lambda^i(f)$ se calculent au moyen de polynômes \emph{universels} à coefficients entiers en les coefficients $c_n(f)=a_n$ de $f$~:
\begin{equation}\tag{1.8}
c_n=P_{i,n}(a_1,\ldots,a_{in}),
\end{equation}
que les relations (1.2) soient vérifiées, et que, pour $i>1$, on ait~:
\begin{equation}\tag{1.9}
\lambda^i(1+at)=1
\end{equation}
($1$ est l'élément nul de l'anneau $K=1+k[[t]]^+$). Les $P_{i,n}$ se calculent encore par des calculs de fonctions symétriques, et l'on constate que $P_{i,n}$ est un polynôme isobare de poids $in$ en les $a_i$. Muni des opérations précédentes, $K$ devient un $\lambda$-anneau, son élément nul est $1$, son élément unité est $1+t$. Notons aussi la relation~:
\begin{equation}\tag{1.7 bis}
(1+at)\circ f=f(at)
\end{equation}

Soit maintenant $K$ un $\lambda$-anneau quelconque. On dit que $K$ est \emph{spécial} si l'homomorphimse additif $\lambda_t$ de $K$ dans $1+K[[t]]^+$ est un $\lambda$-homomorphisme. Cela signifie donc qu'on a les formules~:
\begin{equation}\tag{1.10}
\lambda_t(1)=1+t,\quad
\lambda_t(xy)=\lambda_t(x)\circ\lambda_t(y),\quad
\lambda_t(\lambda^i(x))=\lambda^i(\lambda_t(x),
\end{equation}
ou encore, explicitement~:
\begin{equation}\tag{1.11}
\begin{cases}
\lambda^i(1)=0 & \text{si }i\geq 1\\
\lambda^n(xy)=P_n(\lambda^1x,\ldots,\lambda^n x;\lambda^1y,\ldots,\lambda^n y)\\
\lambda^n(\lambda^i(x))=P_{i,n}(\lambda^1x,\ldots,\lambda^n x),
\end{cases}
\end{equation}
où les $P_n$ et les $P_{i,n}$ sont les polynômes universels envisagés plus haut. Il en résulte en particulier les formules (1.5). (On peut montrer que le $\lambda$-anneau $1+k[[t]]^+$ construit avec un anneau commutatif $k$ est
